\documentclass[11pt]{article}

\usepackage{amsmath,amssymb,amsthm}
\usepackage{graphicx}
\usepackage{hyperref}
\usepackage{booktabs}
\usepackage{enumitem}
\usepackage[margin=1in]{geometry}

\newtheorem{theorem}{Theorem}
\newtheorem{conjecture}{Conjecture}
\newtheorem{remark}{Remark}

\title{Hyperbolic Polytope Volume as an LLM Benchmark Suite:\\
Three Tasks of Increasing Freedom}

\author{Igor Rivin}

\date{\today}

\begin{document}

\maketitle

\begin{abstract}
We propose three complementary benchmark families for evaluating the
mathematical reasoning capabilities of large language models on
constructive geometry: (i) \emph{inscribed reconstruction}, where the
model is given the face lattice of a convex polytope and must produce
vertex coordinates on the unit sphere realizing it; (ii) the
\emph{windowed-volume} task, where the model must construct an ideal
polytope whose hyperbolic volume (or surface area or isoperimetric
ratio) lies inside a narrow tolerance band around a specified target;
and (iii) the \emph{volume lower bound} task, where the model must
report the sharpest lower bound it can on the maximum volume of an
$N$-vertex polytope and back it with explicit vertex coordinates.
The three tasks share a single grading principle---the answer must
be a numerical certificate, evaluatable by a standalone checker---but
they probe different reasoning regimes: pure combinatorial-to-metric
inversion (i), control of a continuous objective (ii), and open-ended
optimization with a falsifiable claim (iii).  We report evaluations of
three frontier models (Claude, GPT-5.4 Pro, Gemini 3.1 Pro DeepThink)
on tasks (i) and (ii), revealing qualitatively different failure
modes; task (iii) is released with full ground truth but the LLM
evaluation is deferred.  As a byproduct of the evaluation we
record a combinatorial invariance phenomenon: the $f$-vector of the
volume-maximizing polytope is the same whether volume is measured
in the Euclidean or hyperbolic sense, across dimensions $d = 3, 4, 5, 6$.
\end{abstract}

% ============================================================================
\section{Introduction}
% ============================================================================

The rapid development of large language models has created a need for
benchmarks that test genuine mathematical reasoning, not just pattern
matching or memorization. An ideal benchmark should:
\begin{enumerate}[label=(\roman*)]
\item have a \emph{verifiable} answer that can be checked automatically;
\item require \emph{construction}, not just classification or selection;
\item be \emph{parameterizable}, so that memorization of specific instances
      does not help;
\item involve \emph{non-trivial mathematics} that is unlikely to appear
      verbatim in training data.
\end{enumerate}

The maximal-volume ideal polytope problem yields three distinct
benchmark families that all satisfy these criteria.  In each case, the
model's response is a list of vertex coordinates on $S^{d-1}$, and the
grader is a standalone Python checker that recomputes the relevant
geometric measure and compares it to a target or record.

% ============================================================================
\section{Three Benchmark Tasks}
% ============================================================================

We summarize the three task families here; per-family results follow
in \S\ref{sec:reconstruction-results}, \S\ref{sec:windowed-results},
and \S\ref{sec:volume-lb-results}.

\paragraph{Task A: inscribed reconstruction.}
Given the $f$-vector, edge list, and facet list of a convex
$d$-polytope, produce $N$ points on $S^{d-1} \subset \mathbb{R}^d$
whose convex hull realizes that combinatorial type.  This is a hard
inverse problem: the model must find a point in a $(dN -
\binom{d+1}{2})$-dimensional configuration space (after factoring out
$\mathrm{O}(d)$) satisfying $O(N^2)$ combinatorial constraints---every
listed facet must appear as a supporting simplex, and no extra
facets may occur.  Two challenge families are released: an
\emph{optimal family} drawn from our volume-maximizing inscribed
polytopes (high automorphism order, often joins or Frobenius
constructions), and a \emph{random family} drawn from random points
on $S^{d-1}$ (generically trivial automorphism group, no exploitable
symmetry).

\paragraph{Task B: windowed volume / surface area / ratio.}
Given integers $N, d$ and a target $T$ with tolerance $\epsilon$,
produce $N$ points $v_1, \ldots, v_N \in S^{d-1}$ such that the
ideal polytope $\operatorname{conv}(v_i) \subset \mathbb{H}^d$ has
the chosen geometric measure (volume, surface area, or isoperimetric
ratio $V/A$) within $[T - \epsilon, T + \epsilon]$.  Targets are
drawn from the $70$th--$95$th percentile of the optimization
landscape, so simple maximization is \emph{insufficient}---the model
must control the geometry to land in a specific window.  The benchmark
includes both hyperbolic and Euclidean variants, with $d = 4$ and
$d = 6$, $N \in \{10, 20, 30\}$.

\paragraph{Task C: volume lower bound (anti-hallucination).}
Given $(d, N)$ and a measure (Euclidean inscribed in the unit ball
of $\mathbb{R}^d$, or hyperbolic ideal in $\mathbb{H}^d$), produce
the sharpest lower bound the model can on the maximum volume,
together with $N$ explicit vertex coordinates that achieve it.  The
explicit-coordinate requirement is the anti-hallucination device: a
model can easily quote a volume number from memory or produce a
plausible-looking one, but it cannot easily fabricate a coordinate
set that the standard convex-hull / Bloch--Wigner / Murakami routine
evaluates to that number.  The benchmark contains $94$ challenges
spanning Euclidean $d \in \{4, 5, 6\}$ and hyperbolic
$d \in \{3, 4, 6\}$.  Hyperbolic $d = 5$ is excluded by design: our
stored $d = 5$ records at $N \in \{12, 14, 20\}$ are known local
rather than global maxima, so grading against them would be
unreliable in that dimension.

\paragraph{What the three tasks share.}
All three answers are graded by a Python checker (see
\S\ref{sec:verification-pipeline}) that returns exit code $0$ for
PASS and $1$ for FAIL.  All three are parameterizable: the
challenge generators emit fresh $(d, N)$ instances on demand.
The mathematics underlying the verification---hyperbolic volume
formulas, Murakami dilogarithms, Hopf reduction in $d=6$---is
sufficiently specialized that training-data contamination is
unlikely.

\paragraph{What the three tasks probe differently.}
Task A is a pure combinatorial-to-metric inversion problem with a
small number of degrees of freedom and many discrete constraints; it
rewards algebraic recognition (``this is the join of two
pentagons'') as well as numerical constraint satisfaction.  Task B
asks the model to control a continuous objective, a strictly harder
test than maximization---a model that can only maximize cannot pass.
Task C is open-ended optimization with a falsifiable claim: any LLM
that beats our record produces new mathematical knowledge.

% ============================================================================
\section{The Verification Pipeline}
\label{sec:verification-pipeline}
% ============================================================================

\subsection{Volume formulas}

For $d = 4$, the volume of an ideal 4-simplex $[\infty, v_1, v_2, v_3, v_4]$
in the upper half-space model is given by the Haagerup--Munkholm formula:
\[
\operatorname{Vol}_4 = \frac{\pi}{3}
\left(4\pi - \sum_{i<j}\varphi_{ij} - \sum_\ell \arccos\frac{d_\ell}{R}\right),
\]
where $\varphi_{ij}$ are the Euclidean dihedral angles of the tetrahedron
$T = \operatorname{conv}(v_1,\ldots,v_4)$, $R$ is its circumradius, and
$d_\ell$ are signed distances from the circumcenter to the faces.

For $d = 6$, ideal simplices have vanishing vertex links ($w_0 = 0$),
and the Hopf reduction gives:
\[
\operatorname{Vol}_6(\Delta_{\mathrm{ideal}})
= \frac{8\pi^3}{15}\left(\frac{17}{4} - w_4 + \frac{1}{2}w_2\right),
\]
where $w_4 = \frac{1}{2\pi}\sum_{i<j}\theta_{ij}$ involves the 21 dihedral
angles, and $w_2 = \frac{1}{2\pi^2}\sum_{|I|=4}\operatorname{Vol}_{S^3}(L_I)$
involves the 35 spherical tetrahedral link volumes, each computable by
Murakami's dilogarithm formula.

\subsection{The standalone checker}

The checker (\texttt{standalone\_checker.py}) is a self-contained Python
script (no external dependencies beyond \texttt{numpy} and \texttt{scipy})
that verifies:
\begin{enumerate}
\item All $N$ vertices lie on $S^{d-1}$ (tolerance $10^{-6}$).
\item All vertices are on the convex hull (no interior points).
\item The convex hull has the prescribed number of facets (if specified).
\item The computed geometric measure falls within $[T-\epsilon, T+\epsilon]$.
\end{enumerate}
Exit code 0 indicates PASS; exit code 1 indicates FAIL.
The checker supports seven measure types: \texttt{volume} ($\mathbb{H}^4$),
\texttt{surface\_area} ($\mathbb{H}^4$), \texttt{ratio} ($\mathbb{H}^4$),
\texttt{volume\_H6}, \texttt{euclidean\_volume}, \texttt{euclidean\_area},
and \texttt{euclidean\_ratio}.

% ============================================================================
\section{Model Evaluation}
% ============================================================================

We evaluated three frontier models---Claude (with the toolkit),
OpenAI's GPT-5.4 Pro (extended reasoning, ``xhigh'' effort), and
Google's Gemini~3.1~Pro DeepThink---on Tasks A and B.  Task C was
designed after the bulk of this evaluation; its model results are
deferred to future work and discussed in
\S\ref{sec:volume-lb-results}.

\subsection{Inscribed reconstruction (Task A)}
\label{sec:reconstruction-results}

We tested GPT-5.4~Pro and Gemini~3.1~Pro DeepThink on instances of
varying difficulty across the optimal and random families.

\begin{table}[ht]
\centering
\caption{LLM performance on inscribed polytope reconstruction.
$\checkmark$ = correct combinatorial type, all points on $S^3$.
$\times$ = wrong type or timeout.}
\label{tab:llm_results}
\smallskip
\begin{tabular}{r l r r c c}
\toprule
$N$ & Family & $|\mathrm{Aut}|$ & Facets & GPT-5.4 & Gemini 3.1 \\
\midrule
10 & optimal & 200  & 25  & ---          & $\checkmark$ \\
12 & random  & 1    & 36  & $\checkmark$ & $\times$ (crash) \\
20 & optimal & 16   & 72  & $\checkmark$ & $\checkmark$ \\
30 & optimal & 1    & 124 & $\checkmark$ & --- \\
\bottomrule
\end{tabular}
\end{table}

The two models employ qualitatively different strategies, as revealed
by their chains of thought.

\paragraph{Gemini: algebraic recognition.}
For the $N = 10$ optimal polytope, Gemini identified the non-edge
structure as two disjoint 5-cycles, recognized the polytope as the
join $C_5 \oplus C_5$ of two regular pentagons in orthogonal 2-planes,
and wrote down the explicit coordinates
$v_j = (\cos 2\pi j/5,\; \sin 2\pi j/5,\; 0,\; 0)$ and
$w_j = (0,\; 0,\; \cos 2\pi j/5,\; \sin 2\pi j/5)$.
This is mathematically correct and produces an exact inscribed
realization.  However, when faced with the $N = 12$ random instance
(trivial automorphism group, no recognizable structure), the same
model crashed without producing output.

\paragraph{GPT: numerical constraint satisfaction.}
For the $N = 12$ random polytope, GPT's chain of thought reveals a
direct numerical approach: ``I've reformulated the search around
determinant-sign tests on the tetrahedral cliques---enforcing
same-side for the 36 listed facets and failure for the others.''
The model effectively runs a gradient-based optimizer over unit
vectors on $S^3$, iteratively adjusting coordinates until the convex
hull matches the target combinatorial type.  This succeeds even for
asymmetric polytopes with no algebraic shortcuts.

\paragraph{GPT on $N = 20$ and $N = 30$.}
The most striking results are GPT's successes on the larger optimal
polytopes.  At $N = 20$ ($|\mathrm{Aut}| = 16$, 72 facets, $54$
continuous DOF after symmetry reduction), GPT's chain of thought
reveals a multi-pronged strategy: it first checked for known polytope
types (24-cell subsets, 600-cell subsets), analyzed vertex figures,
computed graph automorphisms, searched the
Firsching~\cite{Firsching} archive of inscribed polytope realizations,
and ultimately fell back on numerical constraint satisfaction via
determinant-sign tests on facet cliques.  At $N = 30$ with trivial
automorphism group ($|\mathrm{Aut}| = 1$, 124 facets, $80$ effective
continuous DOF), no symmetry reduction is possible, and GPT's success
demonstrates that its in-context numerical engine can navigate an
$80$-dimensional landscape with $124$ nonlinear constraints with no
algebraic shortcuts.  Our own optimizer requires $200$ random
restarts of L-BFGS-B over the same configuration space to reliably
find the $N = 20$ instance.

\paragraph{Gemini on $N = 20$.}
Gemini also solved the $N = 20$ optimal instance with a hybrid
strategy: after failing to identify the polytope as a known type
(24-cell subset, snub 24-cell, join), it computed the orbit
structure under the automorphism group, identifying four orbits of
sizes $[4, 4, 4, 8]$ and the antipodal pairs within each orbit.
It then parameterized the configuration using orbit representatives,
reducing the effective DOF before running L-BFGS-B with facet
hyperplane constraints.  This symmetry-reduced numerical approach
is more sophisticated than a brute-force search over all $3N - 10$
parameters.

\paragraph{Discussion.}
For highly symmetric instances, Gemini's algebraic recognition (as
on $N = 10$) provides an elegant shortcut.  For asymmetric instances,
GPT's direct numerical search (as on $N = 12$ random) succeeds where
Gemini crashes.  At intermediate difficulty ($N = 20$ optimal,
$|\mathrm{Aut}| = 16$), both models converge via numerical methods,
but with different parameterizations.  The inscribed constraint is
essential: without it, the problem reduces to combinatorial
realization (finding \emph{any} coordinates with the right
combinatorics), which is substantially easier.  Indeed, Gemini
solved the non-inscribed $N = 10$ case using non-cospherical
coordinates (norms 2 and $\sqrt{5}$) before we added the $S^3$
requirement.  The hardest instances should be large $N$ with both
high volume and low symmetry---the $N = 16$ optimal polytope
($|\mathrm{Aut}| = 2$, 53 facets, volume per vertex $0.49$) is a
natural candidate.

\subsection{Windowed volume / surface area / ratio (Task B)}
\label{sec:windowed-results}

\subsubsection{Summary of results}

\begin{table}[ht]
\centering
\caption{Model performance on selected benchmark tasks.
Bold indicates the best answer. A dash indicates no numerical answer was produced.}
\label{tab:model-comparison}
\smallskip
\begin{tabular}{l r r r r}
\toprule
Task & Best known & Claude & GPT-5.4 Pro & Gemini 3.1 \\
\midrule
$\mathbb{H}^4$ vol, $N=30$ & 22.262 & \textbf{22.262} & 22.107 & --- \\
$\mathbb{H}^4$ vol, 20 simpl.\ facets & 2.3139 & \textbf{2.3139} & 2.3139 & 1.61$^*$ \\
Eucl.\ vol, $N=30$ on $S^3$ & 2.4909 & \textbf{2.4909} & 2.4839 & --- \\
Eucl.\ vol, 20 facets on $S^3$ & 0.7925 & \textbf{0.7925} & \textbf{0.7925} & \textbf{0.7925}$^\dagger$ \\
Eucl.\ area, 20 facets on $S^3$ & 5.9544 & \textbf{5.9544} & \textbf{5.9544} & \textbf{5.9544}$^\dagger$ \\
Eucl.\ area, 17 facets on $S^3$ & 5.036 & 5.002 & \textbf{5.036} & --- \\
$\mathbb{H}^6$ vol, $N=16$ & 0.5406 & \textbf{0.5406} & 0.480$^\ddagger$ & $0.041^{*}$ \\
\bottomrule
\end{tabular}
\smallskip

{\footnotesize $^*$Gemini's answer is formally correct but useless
(off by a factor of 13 for $\mathbb{H}^6$; stacked-simplex bound for 20 facets).\\
$^\dagger$Gemini provided the correct algebraic formula but no vertex coordinates,
and falsely claimed numerical optimization confirmation.\\
$^\ddagger$GPT-5.4 Pro restricted to centrally symmetric polytopes ($\pm v_i$) and
used QMC volume evaluation; removing the symmetry constraint and using the
exact Hopf reduction formula recovers the full 0.5406.}
\end{table}

\subsubsection{Qualitative comparison}

The three models exhibit qualitatively different strengths and failure modes.

\paragraph{GPT-5.4 Pro: persistent numerical constructor.}
Given a problem, GPT-5.4 Pro typically
produces explicit coordinates, verifies them by multiple methods
(e.g., star triangulations from all vertices), and states honest caveats
(``I have not proved this is optimal''). It identified the free sum
$C_5 \oplus I \oplus I$ as the Euclidean volume maximizer for 20
simplicial facets and derived the exact formula
$V = \frac{5}{24}\sqrt{10+2\sqrt{5}}$.
Its main weakness is occasional misinterpretation of the problem
(e.g., reading $\mathbb{H}^6$ as the 6-cube in one instance, despite
correctly interpreting $\mathbb{H}^4$ in a dozen previous conversations),
which is easily corrected by re-prompting.
On the 17-facet surface area problem, GPT-5.4 Pro found a configuration
(area 5.036) that our 1000-restart optimizer could not beat.

\paragraph{Gemini 3.1 Pro DeepThink: strong algebraic reasoner, unreliable numerical claims.}
Gemini independently derived the exact formula
$A = \frac{5}{3}\sqrt{15-\sqrt{5}}$ for the 20-facet surface area,
and proposed creative constructions (the 10-10 duoprism on the Clifford
torus for 20 non-simplicial facets). However, it claimed that
``random-initialization geometric optimization confirms'' a numerical
value of $12.4417$ for the $N=30$ surface area; when asked to produce
the optimization code, it spent over an hour and delivered code that
only reaches $12.423$. When confronted about the discrepancy, it entered
a computation that ran for over 16 hours with no output.
Its $\mathbb{H}^6$ answer ($4v_6 < M < 220v_6$) is a correct but
useless interval spanning a factor of 55. On a separate occasion,
Gemini produced a long, beautifully structured essay that
misunderstood the optimization direction entirely (minimizing the
supremum instead of maximizing the volume).

\paragraph{Claude (with toolkit): brute-force optimizer with certified answers.}
Using multi-start L-BFGS-B
with 500+ restarts and a standalone verification pipeline, Claude's toolkit
provides the strongest numerical results on most tasks but lacks the
algebraic insight that GPT-5.4 Pro and Gemini occasionally provide. The hybrid
approach---using Euclidean warm starts followed by hyperbolic polishing---was
discovered during this evaluation and reduces optimization time by 38\%.
Several code improvements were adopted from the other models' suggestions:
the Cartesian parameterization with norm projection (from Gemini 3.1's
basin-hopping code) and the Powell optimizer for facet-constrained problems
(from GPT-5.4 Pro's code).

\subsubsection{The ``windowed'' benchmark format}

A key design choice is the \emph{windowed} format: the target value $T$
is chosen from the 70th--95th percentile of the optimization landscape,
not the global maximum. This prevents the model from simply maximizing
the objective (which would give the same answer for every instance) and
instead requires it to \emph{control} the geometry to hit a specific window.
This is substantially harder: a model that can only maximize cannot pass
a windowed test with a target below the maximum.

We generated 66 windowed benchmark cases (33 hyperbolic, 33 Euclidean)
covering volumes, surface areas, and isoperimetric ratios for
$N = 10, 20, 30$ vertices, including facet-count-constrained variants.

\subsubsection{Lessons for LLM training}

The evaluation suggests that current training procedures produce
characteristically different failure modes:
\begin{itemize}
\item \textbf{GPT-5.4 Pro} appears trained to minimize type~II errors
      (failure to produce useful output), resulting in persistent
      numerical construction with honest uncertainty quantification.
\item \textbf{Gemini} appears trained to minimize type~I errors
      (false specific claims), resulting in a tendency to produce
      formulas without evaluating them, and to fabricate numerical
      confirmation when pressured for specific answers.
\item The ability to produce \emph{verifiable certificates} (code
      plus vertex coordinates) separates genuine computation from
      confabulation. Our checker enforces this distinction automatically.
\end{itemize}

\subsection{Volume lower bound (Task C)}
\label{sec:volume-lb-results}

The volume-lower-bound benchmark was designed after the bulk of the
windowed-volume evaluation above, principally as an
anti-hallucination redesign: the windowed benchmark grades the
submitted geometry against a target window and so a model that
produces a plausible-looking number but no coordinates fails
mechanically; Task C makes the requirement explicit by demanding
the coordinates as part of the answer and grading the LLM's claim
against what its own coordinates compute to.

\paragraph{Grading.}
The Task C grader is a single Python script that recomputes the
volume of the submitted polytope under the appropriate measure
(\texttt{ConvexHull} for Euclidean, the closed-form Bloch--Wigner /
Haagerup--Munkholm / Murakami routines for hyperbolic), checks
consistency between the LLM's claimed bound and the recomputed
value within $10^{-3}$, and compares the recomputed value to our
internal record:
\begin{itemize}
\item \textbf{strong:} computed $\geq$ record (the model has matched
      or beaten our published bound; new mathematical claim);
\item \textbf{pass:} computed $\geq 0.95 \times$ record;
\item \textbf{fail:} valid polytope but $< 0.95 \times$ record;
\item \textbf{invalid:} no coordinates, wrong dimension, or
      non-positive volume.
\end{itemize}

\paragraph{Coverage.}
The benchmark contains $94$ challenges across five regimes:

\begin{center}
\begin{tabular}{l c c}
\toprule
Measure & Dimension & $N$ values (count) \\
\midrule
Euclidean inscribed & $d = 4$ & $5$--$24$ ($20$) \\
Euclidean inscribed & $d = 5$ & $6$--$20$ ($15$) \\
Euclidean inscribed & $d = 6$ & $7$--$24$ ($18$) \\
Hyperbolic ideal    & $d = 3$ & $4$--$300$ ($17$, sparse at large $N$) \\
Hyperbolic ideal    & $d = 4$ & $5$--$30$ ($14$) \\
Hyperbolic ideal    & $d = 6$ & $7$--$24$ ($10$) \\
\bottomrule
\end{tabular}
\end{center}

Hyperbolic $d = 5$ is excluded by design: our stored configurations
at $N \in \{12, 14, 20\}$ are local rather than global maxima
of the Euclidean (and hence hyperbolic) objective in $d = 5$, so
grading against them would be unreliable for those instances and
artificially easy at $N \in \{7, 8, 9, 10, 16\}$ where the records
do match the global maxima.

\paragraph{Status.}
The benchmark files (\texttt{benchmark.json},
\texttt{ground\_truth.json}, README) and grader
(\texttt{check\_volume\_lb\_response.py}) are released alongside
this note.  Frontier-model evaluation on Task C is deferred to
future work; the design and ground-truth pipeline are reported
here so the benchmark can be used by other groups.

% ============================================================================
\section{Mathematical Byproducts}
% ============================================================================

The LLM evaluation process, combined with systematic computational
exploration, led to several mathematical discoveries.

\subsection{Combinatorial invariance}
The $f$-vector of the volume-maximizing ideal polytope is the same
whether volume is measured in the Euclidean or hyperbolic sense.
This was verified in $d = 3$ ($N \leq 50$), $d = 4$ ($N \leq 30$),
and $d = 6$ ($N \leq 24$), with no exceptions. The invariance holds
both within the regime where optimal polytopes are product polytopes
and beyond it, suggesting a universal phenomenon.

\subsection{Identification of optimal polytopes}
For small $N$ in even dimension $d = 2k$, the volume-optimal polytope
is the join of $k$ regular polygons in orthogonal 2-planes:
$C_{p_1} \oplus \cdots \oplus C_{p_k}$. This is optimal up to
$N \approx 2d + 2$, beyond which asymmetric configurations take over.
Notable cases include the cross-polytope ($N = 2d$) and the simplex ($N = d+1$).
Volume-optimal polytopes are not cyclic polytopes and are in fact far
from neighborly: at $N = 30$ in $d = 4$, the optimal polytope has
only 123 facets compared to 756 for the cyclic polytope $C(30,4)$.

\subsection{The volume plateau}
At $N = 9$ in $d = 4$ (equivalently, 20 simplicial facets), the
hyperbolic volume maximum is achieved on a manifold of dimension $\geq 2$,
not an isolated point. The Euclidean maximizer (the free sum
$C_5 \oplus I \oplus I$, with exact volume $\frac{5}{24}\sqrt{10+2\sqrt{5}}$)
and the hyperbolic maximizer are metrically distinct configurations with
identical hyperbolic volume to 13 significant digits. For prime facet
counts (e.g., 17), the maximum is isolated.

\subsection{Volume--volume correlation}
Across the optimization landscape, the Euclidean and hyperbolic volumes
are correlated at $r > 0.995$, with log-log elasticity approaching 2
at $N = 30$ in $d = 4$. This near-proportionality is surprising given
that the Euclidean volume is polynomial in the vertex coordinates while
the hyperbolic volume involves the transcendental Haagerup--Munkholm formula.
For random configurations (as opposed to near-optimal ones), the correlation
is lower ($r \approx 0.86$ in $d = 6$) but still substantial, and holds
even within a fixed combinatorial type ($r \approx 0.88$ within a single
facet-count bucket).

\subsection{Distribution of random polytope volumes}

For $N$ points sampled uniformly on $S^{d-1}$, both the Euclidean
hull volume and the ideal $\mathbb{H}^d$ polytope volume are random
variables whose distribution depends on $N$ and $d$.
In $d = 6$, we find a transition in the distribution shape as $N$ grows:

\begin{center}
\begin{tabular}{r r r r l}
\toprule
$N$ & Skewness & Normal $p$ & Beta $p$ & Best fit \\
\midrule
10 & $+0.78$ & $<0.001$ & 0.985 & Beta$(3.9, 43.7)$ \\
12 & $+0.65$ & 0.005 & 0.766 & Beta$(6.9, 120.7)$ \\
16 & $+0.23$ & 0.088 & $<0.001$ & Borderline \\
20 & $+0.02$ & \textbf{0.832} & $<0.001$ & Normal \\
30 & $-0.09$ & 0.281 & $<0.001$ & Normal \\
50 & $-0.19$ & 0.034 & $<0.001$ & Normal (slight left skew) \\
\bottomrule
\end{tabular}
\end{center}

For small $N$ ($\leq 12$), the volume distribution is well-described by
a Beta distribution with parameters scaling linearly in $N$, consistent
with earlier observations in $d = 3$ and $d = 4$.
For large $N$ ($\geq 20$), the distribution converges to Normal---as
expected from a CLT-type argument, since the polytope volume becomes a
sum of many weakly-dependent simplex volumes in the Delaunay triangulation.
The crossover occurs at $N \approx 16\text{--}20$ in $d = 6$ (where the
skewness crosses zero), and is accompanied by a subtle transition from
right skew to slight left skew.

The within-bucket distributions (conditioning on a fixed facet count) show
essentially the same shape and correlation structure, confirming that the
distributional properties are genuinely geometric rather than combinatorial
artifacts.

% ============================================================================
\section{Conclusions}
% ============================================================================

The three benchmark families above expose qualitative differences
between frontier LLMs that traditional math benchmarks miss. The
ability to produce \emph{verifiable certificates}---explicit
coordinates that pass an independent numerical checker---is a
fundamentally different capability from producing plausible-sounding
mathematical text.  The three tasks probe different capabilities:
inscribed reconstruction (Task A) is a constraint-satisfaction
problem rewarding both algebraic recognition and numerical search;
the windowed-volume task (Task B) requires controlling a continuous
objective rather than just maximizing it; and the volume-lower-bound
task (Task C) is anti-hallucination by design---a numerical claim
that does not match the model's own coordinates is automatically
rejected.

All three benchmarks are parameterizable (any $N$, $d$, and either
target window or measure), resistant to memorization (windowed
targets from a continuous landscape; record values withheld in
Task C), and fully automatically gradable (standalone checkers with
exit codes).  The mathematical content---hyperbolic volume,
Murakami dilogarithms, Hopf reduction, Bloch--Wigner---is
sufficiently specialized that training-data contamination is
unlikely.

As a byproduct, the evaluation process itself generated new mathematical
results (combinatorial invariance, identification of optimal polytopes,
volume plateau), demonstrating that the ``mixture of experts'' approach
---combining human mathematical intuition with LLM algebraic insight
and automated numerical optimization---can be more productive than
any single participant acting alone.

% ============================================================================
% References
% ============================================================================

\begin{thebibliography}{99}

\bibitem{HM}
U.~Haagerup and H.~J.~Munkholm,
``Simplices of maximal volume in hyperbolic $n$-space,''
\emph{Acta Math.} \textbf{147} (1981), 1--11.

\bibitem{Murakami}
J.~Murakami,
``Volume formulas for a spherical tetrahedron,''
\emph{Osaka J. Math.} \textbf{49} (2012), 631--638.

\bibitem{RivinBasmajian}
I.~Rivin,
``On Basmajian's identities, and other stories,''
arXiv:1404.1583, 2014.

\bibitem{H56note}
I.~Rivin,
``Notes on hyperbolic 5- and 6-simplex volumes,''
companion note, 2026.

\bibitem{Firsching}
M.~Firsching,
``The complete enumeration of 4-polytopes and 3-spheres with nine vertices,''
\emph{Israel J. Math.} (2020).

\end{thebibliography}

\end{document}
